\documentclass[11pt,a4paper]{article}
\usepackage[utf8]{inputenc}
\usepackage[T1]{fontenc}
\usepackage[portuguese]{babel}
\usepackage[margin=2.5cm]{geometry}
\usepackage{booktabs}
\usepackage{lmodern}

\title{PCS5916 Computação de Alto Desempenho e Aplicações\\[4pt]\large Relatório dos Exercícios da Aula 03}
\author{André Moreno Goveia\\13682785}
\date{4 de outubro de 2026}

\begin{document}
\sloppy
\maketitle

\section{Ambiente de testes}

Os testes rodaram em um processador de 8 núcleos e 16 threads. Compilei os programas com \texttt{gcc -fopenmp} usando \texttt{-O0} e \texttt{-O2} e rodei cada versão 10 vezes. Os tempos das tabelas são médias e medem só o trecho de cálculo, sem a inicialização dos dados. O speedup é o tempo sequencial dividido pelo tempo paralelo. Em todos os casos as versões paralelas deram os mesmos resultados das sequenciais.

\section{Exercício 1: estatísticas de um vetor}

O programa calcula a soma, o maior e o menor elemento e a quantidade de pares de um vetor com 100 milhões de inteiros positivos. Na versão sequencial, um único laço percorre o vetor e atualiza as quatro estatísticas ao mesmo tempo.

\subsection{Versão com parallel for (ex1\_for.c)}

A ideia aqui é dividir o vetor. Com \texttt{\#pragma omp parallel for num\_threads(4)}, cada thread fica com um quarto dos elementos e calcula as quatro estatísticas na sua parte. Como todas as threads atualizam as mesmas variáveis, usei \texttt{reduction(+:soma, pares)}, \texttt{reduction(max:maior)} e \texttt{reduction(min:menor)}. Assim cada thread trabalha com uma cópia própria e o OpenMP junta os resultados no fim, sem precisar de região crítica.

\subsection{Versão com parallel sections (ex1\_sections.c)}

Nesta versão a divisão é por tarefa. Criei quatro seções, uma para cada estatística, e cada thread percorre o vetor inteiro calculando só a sua. Cada seção escreve em uma variável diferente, então não há condição de corrida.

Na primeira tentativa as seções escreviam direto nas variáveis \texttt{soma}, \texttt{maior}, \texttt{menor} e \texttt{pares}. Essas variáveis ficam lado a lado na memória e caem na mesma linha de cache, e por isso cada escrita de uma thread invalidava a cache das outras. Esse falso compartilhamento deixou o programa com cerca de 0,82 s em \texttt{-O2}, mais de dez vezes mais lento que o sequencial. Para resolver, cada seção passou a acumular em uma variável local e só copia o resultado no final.

\subsection{Resultados}

\begin{table}[h]
\centering
\begin{tabular}{lcccc}
\toprule
 & \multicolumn{2}{c}{\texttt{-O0}} & \multicolumn{2}{c}{\texttt{-O2}} \\
Versão & Tempo (s) & Speedup & Tempo (s) & Speedup \\
\midrule
Sequencial (ex1\_seq.c) & 1,2089 & 1,00 & 0,0738 & 1,00 \\
parallel for (ex1\_for.c) & 0,3275 & 3,69 & 0,0779 & 0,95 \\
parallel sections (ex1\_sections.c) & 1,0544 & 1,15 & 0,0635 & 1,16 \\
\bottomrule
\end{tabular}
\caption{Tempos médios de 10 execuções do exercício 1 com 4 threads.}
\end{table}

Sem otimização, o \texttt{parallel for} chegou a 3,69, bem perto do ideal de 4, já que cada thread faz um quarto do trabalho e juntar as reduções custa quase nada. O \texttt{parallel sections} ficou em 1,15. Faz sentido, porque cada thread lê o vetor inteiro e o tempo total acaba sendo o da seção mais lenta. Além disso, essa abordagem não escala, porque com só quatro estatísticas uma quinta thread não teria o que fazer.

Com \texttt{-O2} o cenário muda. O compilador vetoriza o laço e a versão sequencial lê os 400 MB do vetor em cerca de 74 ms. Nesse ponto o gargalo passa a ser a memória e não o processamento, então colocar mais threads não ajuda e o \texttt{parallel for} fica até um pouco mais lento por causa do custo de criar as threads. O \texttt{parallel sections} ganhou 16\%, provavelmente porque cada seção tem um laço mais simples, que o compilador otimiza melhor. No fim, esse problema faz muito pouca conta por elemento lido, e o paralelismo só vale a pena quando o cálculo pesa mais que o acesso à memória.

\section{Exercício 2: multiplicação de matrizes}

Parti do \texttt{mat\_mul.c} da aula 02, que multiplica duas matrizes $1000 \times 1000$ com três laços aninhados. Cada linha de $C$ depende só de uma linha de $A$ e da matriz $B$, então linhas diferentes podem ser calculadas em paralelo sem nenhuma sincronização.

\subsection{Versão com parallel for (ex2\_for.c)}

Coloquei \texttt{\#pragma omp parallel for num\_threads(4) private(j, k)} no laço externo, e cada thread calcula 250 linhas de $C$. O \texttt{private(j, k)} é necessário porque no código original essas variáveis são declaradas no início do \texttt{main} e seriam compartilhadas entre as threads. O \texttt{i} já é privado por ser o índice do laço paralelizado.

\subsection{Versão com parallel sections (ex2\_sections.c)}

Movi os três laços para uma função \texttt{multiplica(inicio, fim)} que calcula um intervalo de linhas de $C$. Cada uma das quatro seções chama essa função para um bloco de 250 linhas. Na prática é a mesma divisão que o \texttt{parallel for} faz, só que feita à mão. Como todas as linhas têm o mesmo custo, a carga fica equilibrada.

\subsection{Resultados}

\begin{table}[h]
\centering
\begin{tabular}{lcccc}
\toprule
 & \multicolumn{2}{c}{\texttt{-O0}} & \multicolumn{2}{c}{\texttt{-O2}} \\
Versão & Tempo (s) & Speedup & Tempo (s) & Speedup \\
\midrule
Sequencial (mat\_mul.c) & 6,4245 & 1,00 & 0,6315 & 1,00 \\
parallel for (ex2\_for.c) & 2,1242 & 3,02 & 0,2644 & 2,39 \\
parallel sections (ex2\_sections.c) & 2,0383 & 3,15 & 0,2532 & 2,49 \\
\bottomrule
\end{tabular}
\caption{Tempos médios de 10 execuções do exercício 2 com 4 threads.}
\end{table}

As duas versões tiveram speedups parecidos, 3,02 e 3,15 sem otimização, o que era esperado já que dividem o trabalho do mesmo jeito. A diferença entre elas fica dentro da variação normal entre execuções. O speedup não chega a 4 principalmente por causa do acesso a $B[k][j]$, que percorre uma coluna e pula 4000 bytes a cada leitura, gerando muitas faltas de cache. Com quatro threads essas faltas disputam a mesma cache L3 e a mesma memória. O processador também baixa a frequência quando mais núcleos estão ativos. Com \texttt{-O2} o cálculo fica cerca de dez vezes mais rápido, a memória passa a pesar mais e o speedup cai para cerca de 2,4. Mesmo assim o ganho é bem maior que no exercício 1, porque aqui são $10^9$ multiplicações sobre apenas 12 MB de dados.

\section{Exercício 3: jantar dos filósofos}

Cinco filósofos sentam em volta de uma mesa com um garfo entre cada par, e cada um precisa dos dois garfos ao seu lado para comer. Nas duas versões, \texttt{\#pragma omp parallel num\_threads(5)} cria uma thread por filósofo. O filósofo $i$ usa o garfo $i$ à esquerda e o garfo $(i+1) \bmod 5$ à direita. Cada um pensa por até 1 ms, pega os garfos, come por 0,5 ms e devolve os garfos. A simulação dura 3 s e, no final, o programa mostra quantas vezes cada filósofo comeu e há quanto tempo estava sem comer. Coloquei também uma pausa de 1 ms entre pegar o primeiro e o segundo garfo. Ela não muda a lógica, mas deixa os problemas de sincronização bem mais fáceis de observar.

\subsection{Versão padrão (ex3\_padrao.c)}

Cada garfo é um inteiro do vetor \texttt{garfos}, que vale 1 quando está em uso. O filósofo fica em um laço esperando o garfo esquerdo ficar livre, marca como ocupado, faz o mesmo com o direito, come e solta os dois. O vetor é \texttt{volatile} para o compilador não guardar o valor em registrador, o que transformaria a espera em um laço infinito. Os laços de espera também olham o relógio, mas só para o programa conseguir terminar depois dos 3 s.

Essa versão tem três problemas. Primeiro, existe uma condição de corrida, porque testar o garfo e marcá-lo são dois passos separados. Dois vizinhos podem ver o mesmo garfo livre ao mesmo tempo e os dois pegarem.

Segundo, pode acontecer deadlock. Se todos pegam o garfo esquerdo ao mesmo tempo, cada um fica esperando o garfo que o vizinho da direita está segurando, e ninguém sai do lugar. Com a pausa de 1 ms isso aconteceu logo na primeira rodada em todas as execuções, e todos os filósofos terminaram sem comer nenhuma vez. Tirando a pausa, o deadlock continuou aparecendo, só que em momentos aleatórios, depois de 115 a 867 refeições por filósofo.

Terceiro, pode haver starvation. Nada garante que um filósofo consiga os dois garfos, então um que perde a disputa várias vezes seguidas pode ficar sem comer enquanto os outros comem. E como a espera é ativa, cada filósofo esperando ocupa um núcleo inteiro à toa.

\subsection{Versão com locks (ex3\_locks.c)}

Aqui cada garfo é um \texttt{omp\_lock\_t}. O filósofo pega o garfo com \texttt{omp\_set\_lock}, que bloqueia a thread até o lock ficar livre, e devolve com \texttt{omp\_unset\_lock}. O OpenMP garante que só uma thread segura cada lock por vez, o que resolve a condição de corrida e acaba com a espera ativa.

Só trocar por locks não resolve o deadlock, porque se todos pegarem o garfo esquerdo primeiro o ciclo de espera continua igual. Para quebrar o ciclo, cada filósofo pega primeiro o garfo de menor número. Os filósofos 0 a 3 continuam começando pela esquerda, mas o filósofo 4 pega o garfo 0 antes do garfo 4. Com isso a espera circular não consegue mais se formar.

Os locks do OpenMP também não garantem ordem de chegada, então um filósofo ainda poderia perder a disputa repetidamente. Para evitar isso, adicionei uma regra simples. Antes de pegar os garfos, o filósofo confere se já comeu mais que algum vizinho e, se sim, espera os vizinhos o alcançarem. Os contadores usam \texttt{\#pragma omp atomic} para leitura e escrita. Quem comeu menos nunca precisa esperar, por isso a mesa sempre avança. Como a espera acontece sem segurar garfo nenhum, ela também não cria um novo deadlock.

Nos testes, mesmo com a pausa de 1 ms, a versão com locks rodou os 3 s sem travar. Cada filósofo comeu entre 859 e 861 vezes e nenhum ficou mais de 2 ms sem comer no final.

\end{document}
